\section{Randomized --- BPPHard}
%%% generated by the blueprint pipeline from TCSlib.Complexity.Randomized.BPPHard
\subsection*{Randomized --- BPPHard: Introduction}

Declarations of the Lean module \texttt{TCSlib.Complexity.Randomized.BPPHard}.

\subsection*{Randomized --- BPPHard: Results}

\begin{definition}[Binary-choice machine code]
\lean{Turing.CodePTM}
\label{Turing.CodePTM}
\leanok
\uses{Turing.NDTM}
A code for a binary-choice Turing machine consists of a natural number $k$ of work
tapes, a natural number $m$ such that the machine has state set $	ext{Fin}(m+1)$, and
the machine itself, whose input alphabet is Boolean and whose transitions are indexed by
a choice bit.
\end{definition}

\begin{definition}[Finite-state machine conversion]
\lean{Turing.CodePTM.toFinNDTM}
\label{Turing.CodePTM.toFinNDTM}
\leanok
\uses{Turing.CodePTM, Turing.FinNDTM}
Given a binary-choice Turing machine code with $k$ work tapes and state set
$	ext{Fin}(m+1)$, its associated finite-state machine has $k$ work tapes, that state
set, and the same underlying machine.
\end{definition}

\begin{definition}[Padded machine coding]
\lean{Turing.PTMCode}
\label{Turing.PTMCode}
\leanok
\uses{Turing.CodePTM}
A padded machine coding consists of an encoding of every binary-choice Turing machine
with finitely many work tapes and states as a finite binary string, a decoding map from
every finite binary string to such a machine, and the property that appending any finite
number of $1$ bits to a machine's code still decodes to that machine.
\end{definition}

\begin{theorem}[Decoding an encoded machine]
\lean{Turing.PTMCode.decode_encode}
\label{Turing.PTMCode.decode_encode}
\leanok
\uses{Turing.CodePTM, Turing.PTMCode}
\difficulty{1}
Let $c$ be a machine code with an encoder and decoder satisfying the padding property
that decoding the encoding of any machine followed by any number of true bits returns
that machine. For every machine $M$, decoding its encoding returns $M$.
\end{theorem}

\begin{definition}[Tuple of machine actions]
\lean{Turing.actionTuple}
\label{Turing.actionTuple}
\leanok
\uses{Turing.Action}
For an action of a machine with $k$ work tapes, Boolean tape symbols, and state set $S$,
its action tuple is the tuple consisting of the input-head movement, the function
assigning each work tape its write instruction and head movement, the optional output
bit, and the optional successor state.
\end{definition}

\begin{theorem}[Injectivity of the action record]
\lean{Turing.actionTuple_injective}
\label{Turing.actionTuple_injective}
\leanok
\uses{Turing.actionTuple}
\difficulty{4}
For any natural number $k$ and any type $S$, the map sending an action with $k$ work
tapes, binary symbols, and states in $S$ to the tuple of its input-head move, work-tape
data, output symbol, and successor state is injective.
\end{theorem}

\begin{definition}[Machine data as a dependent tuple]
\lean{Turing.codeSigma}
\label{Turing.codeSigma}
\leanok
\uses{Turing.Action, Turing.CodePTM}
For a binary-choice Turing machine code with $k$ work tapes and one more than its stated
number of states, the associated dependent tuple records the pair of natural numbers
$(k,	ext{number of states minus one})$, an initial state among the corresponding states,
and the two transition functions specifying the machine's actions.
\end{definition}

\begin{theorem}[Injectivity of machine encoding]
\lean{Turing.codeSigma_injective}
\label{Turing.codeSigma_injective}
\leanok
\uses{Turing.codeSigma}
\difficulty{4}
A binary-choice Turing machine with $k$ work tapes and state set $\{0,\ldots,s\}$ is
determined by the tuple consisting of $k$, $s$, its initial state, and its two
transition functions. In particular, the map sending such a machine to this tuple is
injective.
\end{theorem}

\begin{theorem}[Existence of a machine code]
\lean{Turing.PTMCode.nonempty}
\label{Turing.PTMCode.nonempty}
\leanok
\uses{Turing.CodePTM, Turing.PTMCode}
\difficulty{4}
There exists a probabilistic Turing machine code, consisting of an encoding and decoding
of machine descriptions such that decoding an encoding followed by any number of true
bits recovers the original machine.
\end{theorem}

\begin{definition}[Acceptance probability within a time bound]
\lean{Turing.FinNDTM.acceptWithinProb}
\label{Turing.FinNDTM.acceptWithinProb}
\leanok
\uses{Randomized.randProb, Turing.FinNDTM, Turing.FinNDTM.outputWith, Turing.NDTM.initCfg, Turing.NDTM.runWith}
For a finite nondeterministic Turing machine $N$ over the Boolean alphabet, a Boolean
input $x$, and a time bound $t$, the acceptance probability within the time bound is a
rational number assigned to $(N,x,t)$.
\end{definition}

\begin{theorem}[Coding a Binary-Choice Machine]
\lean{Turing.FinNDTM.exists_codePTM}
\label{Turing.FinNDTM.exists_codePTM}
\leanok
\uses{Turing.Cfg.mapState, Turing.CodePTM, Turing.CodePTM.toFinNDTM, Turing.FinNDTM, Turing.FinNDTM.outputWith, Turing.NDTM.initCfg, Turing.NDTM.relabelState, Turing.NDTM.relabelState_runWith, Turing.NDTM.runWith}
\difficulty{4}
For every finite-state nondeterministic Turing machine $N$ over the binary alphabet,
there is a machine $M$ with a finite number of work tapes, states indexed by an initial
segment of the natural numbers, and binary-choice transitions such that, for every
binary input $x$ and every choice word $w$, the run of $M$ halts if and only if the run
of $N$ halts, and both runs produce exactly the same output.
\end{theorem}

\begin{theorem}[Acceptance probability equals output probability]
\lean{Turing.FinNDTM.acceptWithinProb_eq_outputProb}
\label{Turing.FinNDTM.acceptWithinProb_eq_outputProb}
\leanok
\uses{Randomized.randProb_congr, Turing.FinNDTM, Turing.FinNDTM.acceptWithinProb, Turing.FinNDTM.outputProb, Turing.FinNDTM.outputProb_eq_of_haltsWithin, Turing.NDTM.HaltsWithin, Turing.NDTM.HaltsWithin.mono}
\difficulty{4}
Let $N$ be a finite nondeterministic Turing machine over Boolean symbols, let $x$ be a
Boolean input, and let $T,t\in\mathbb{N}$. If every run on $x$ given a choice word of
length $T$ halts, and $T\le t$, then the probability that $N$ accepts within $t$ steps
equals the probability that a uniformly chosen Boolean choice word of length $T$
produces output $[\mathrm{true}]$.
\end{theorem}

\begin{definition}[Probabilistic acceptance language]
\lean{Randomized.ProbAccept}
\label{Randomized.ProbAccept}
\leanok
\uses{Turing.CodePTM.toFinNDTM, Turing.FinNDTM.acceptWithinProb, Turing.PTMCode, Turing.pairEncode}
Given a padded machine coding, its probabilistic acceptance language consists of binary
strings $y$ for which there exist binary strings $\alpha$ and $x$, and a natural number
$t$, such that $y$ is the prescribed pair encoding of $\alpha$ and the pair encoding of
$x$ followed by $t$ copies of $1$, and the machine decoded from $\alpha$ accepts $x$
within time $t$ with probability at least $2/3$.
\end{definition}

\begin{definition}[BPP-hard language]
\lean{Randomized.BPPHard}
\label{Randomized.BPPHard}
\leanok
\uses{Randomized.BPP_PTM}
A language $L$ over the Boolean alphabet is BPP-hard if every language $L'$ in the class
of languages decided by probabilistic polynomial-time Turing machines is polynomial-time
many-one reducible to $L$.
\end{definition}

\begin{theorem}[A Power Bound by a Shifted Power]
\lean{Randomized.pow_add_one_le_two_mul}
\label{Randomized.pow_add_one_le_two_mul}
\leanok
\difficulty{3}
For all natural numbers $n,e$, one has $n^e+1\le 2(n+1)^e$.
\end{theorem}

\begin{theorem}[Hardness of a Probabilistic Machine Acceptance Language]
\lean{Randomized.probAccept_bppHard}
\label{Randomized.probAccept_bppHard}
\leanok
\uses{Complexity.PolyTimeComputable.pairEncode, Complexity.polyTimeComputable_const, Complexity.polyTimeComputable_id, Complexity.polyTimeComputable_polyUnary, Randomized.BPPHard, Randomized.ProbAccept, Randomized.pow_add_one_le_two_mul, Randomized.randProb_congr, Randomized.randProb_mono, Randomized.randProb_not, Turing.CodePTM.toFinNDTM, Turing.FinNDTM.acceptWithinProb, Turing.FinNDTM.acceptWithinProb_eq_outputProb, Turing.FinNDTM.exists_codePTM, Turing.FinNDTM.outputProb, Turing.MultiTapeTM.indicator, Turing.PTMCode, Turing.PTMCode.decode_encode, Turing.pairEncode, Turing.pairEncode_injective}
\difficulty{6}
For every code $c$ of a binary-choice probabilistic Turing machine, the language whose
strings encode triples $(\alpha,x,t)$ as $\langle\alpha,\langle x,1^t\rangle\rangle$ and
for which the machine decoded from $\alpha$ accepts $x$ within $t$ steps with
probability at least $2/3$ is hard for $\mathsf{BPP}$ under polynomial-time many-one
reductions. Here $\mathsf{BPP}$ is the class of languages decidable by a finite
probabilistic Turing machine in polynomial time with acceptance probability at least
$2/3$ on members and at most $1/3$ on nonmembers.
\end{theorem}
